% ============================================================
%  1. INTRODUCTION
% ============================================================
\chapter{Introduction (Work in progress (EGON))}\label{sec:introduction}
The Phillips curve, originally proposed by \cite{phillips1958} and later 
formalized by \cite{samuelson1960}, postulates an inverse relationship between 
inflation and unemployment. Despite its theoretical appeal, the empirical 
stability of this relationship has been widely debated. In particular, the 
curve appeared to break down during the stagflation of the 1970s, and more 
recently, the post-COVID surge in inflation has reignited the discussion of 
whether unemployment retains meaningful predictive power over inflation dynamics 
\parencite{stock2019}.

This paper investigates whether the Phillips curve relationship carries 
forecasting value for Danish inflation. Specifically, we ask: 
\begin{center}
    \textit{Can we produce more accurate forecasts of Danish inflation by incorporating unemployment data and the broader European macroeconomic environment, compared to a purely univariate benchmark?}
\end{center}
To answer this question, we employ two forecasting models. First, an 
ARIMA model serves as a univariate baseline, forecasting Danish inflation using 
only its own history. Second, a Vector Autoregression (VAR) model incorporates 
Danish unemployment, Danish inflation, and Eurozone inflation jointly, allowing 
us to assess whether cross-variable dynamics -- including potential spillovers 
from the European economy -- improve forecast accuracy. Variable selection for 
the VAR model is guided by a General-to-Specific (GETS) approach, following 
\cite{doornik2014}. Forecast accuracy is evaluated using standard measures 
including RMSE and MAE, and models are compared both in-sample and 
out-of-sample.

The remainder of the paper is structured as follows. Section 2 describes the 
data. Section 3 introduces the methodology. Section 4 presents and discusses 
the results. Section 5 concludes.

